\subsection{Cohort Definitions}
This section is verbatim reproduced from the conversation with JD Wright to provide
background on the history of cohorts and their rationale.

\textbf{Cohort 1}: Boxed parasiticides. This was the best starting point
because it represents roughly 40\% of total DUR volume and the
majority of the initial digital volume, while the clinical rules are
relatively standardized and label-driven. The primary checks are
species, weight range, monthly directions, days\textquotesingle{}
supply, duplicate parasite coverage, and known contraindications. This
gave us a high-volume population with relatively low variability for
validating the core RxBuddy and DUR Copilot experience.

\textbf{Cohort 2}: Non-controlled, non-boxed medications. This represents
roughly 29\% of total DUR volume and materially expands coverage, but
the clinical review is more complex. These medications require
generalized dosing evaluation, including strength, dose, frequency,
route, indication, species appropriateness, duplicate therapy,
contraindications, and potential interactions. We separated this
cohort so those broader AI checks could be validated without also
introducing compounding or controlled-substance requirements at the
same time.

\textbf{Cohort 3}: Non-controlled compounded medications. This is a smaller
segment at roughly 4.4\% of DUR volume, but it has distinct clinical
and compliance requirements. Compounds require a documented
compounding reason, formulation-specific directions, and validation
that the prescribed days\textquotesingle{} supply does not exceed the
product\textquotesingle s beyond-use date. BUD rules vary by dosage
form, storage conditions, formulation, and stability-study results.
Keeping compounds separate allowed us to test these specialized
workflows and manage their impact on pharmacist throughput
independently.

\textbf{Cohort 4}: Gabapentin and the high-scrutiny regulatory workflow. This
represents roughly 1.1\% of DUR volume. The implemented Phase 4 scope
focused on uncontrolled gabapentin rather than the full
controlled-substance population, but it was separated because it
introduces site eligibility, pharmacist skill gating, regulatory
indicators, mixed-order handling, and state-specific scrutiny. It also
gave us a contained way to validate infrastructure that will be needed
for broader controlled-substance support later.
